\documentclass[11pt]{article}
\usepackage[margin=1in]{geometry}
\usepackage{amsmath,amssymb,amsthm,mathtools}
\usepackage{microtype}
\usepackage[hidelinks]{hyperref}
\usepackage{enumitem}
\setlength{\emergencystretch}{1em}
\newtheorem{theorem}{Theorem}
\newtheorem{proposition}[theorem]{Proposition}
\newtheorem{lemma}[theorem]{Lemma}
\theoremstyle{remark}
\newtheorem{remark}[theorem]{Remark}
\newcommand{\R}{\mathbb R}
\newcommand{\E}{\mathbb E}
\newcommand{\Prob}{\mathbb P}
\newcommand{\ip}[2]{\langle #1,#2\rangle}
\newcommand{\norm}[1]{\lVert #1\rVert}
\newcommand{\dist}{\operatorname{dist}}
\newcommand{\Risk}{\mathcal R}
\newcommand{\Class}{\mathcal F}
\newcommand{\Pol}{\mathcal P}
\title{Adaptive random reshuffling:\\a matching worst-case lower bound}
\author{}
\date{}

\begin{document}
\maketitle

\begin{abstract}
Consider random reshuffling of a smooth convex finite sum with $n$
components over $K$ epochs. We prove a matching worst-case lower bound
when every update uses a positive scalar step at most $1/L$ and the
output is any scalar convex combination of the actual iterates. Steps
and output weights may depend on the instance, private randomness, and
even the entire bank of sampled permutations. For every parameter tuple,
one instance of dimension at most $n$ is hard for all such choices.
Its lower bound is a universal constant times
\[
 \min\!\left\{LD^2,\frac{LD^2}{nK}
 +\min\!\left\{\frac{\sigma D}{\sqrt{nK}},
       \left(\frac{L\sigma^2D^4}{nK^2}\right)^{1/3}\right\}\right\}.
\]
The matching upper bound follows from Liu's constant-step,
pre-update-average guarantee and the option of returning the initial
point. The lower uses regular-simplex offsets, a squared-distance cone
potential, and a pathwise interval estimate that pays for every
auxiliary reset. A shortest-maximizing-interval argument controls the
endpoint terms without imposing confinement on the actual trajectory.
\end{abstract}

\section{Model and result}

Fix integers $n\ge2$, $K\ge1$, put $T=nK$, and let $L>0$ and
$D,\sigma\ge0$. The class $\Class(n,L,D,\sigma)$ consists of instances
in arbitrary finite-dimensional Euclidean spaces with differentiable
convex components $f_1,\ldots,f_n$ such that
\begin{equation}\label{eq:class}
 \norm{\nabla f_i(x)-\nabla f_i(y)}\le L\norm{x-y},\qquad
 f=\frac1n\sum_{i=1}^n f_i,\qquad
 x_*\in\arg\min f,
\end{equation}
and specified initial point $x_0$ satisfying
\begin{equation}\label{eq:parameters}
 \norm{x_0-x_*}\le D,\qquad
 \frac1n\sum_{i=1}^n\norm{\nabla f_i(x_*)}^2\le\sigma^2.
\end{equation}
No global bound on gradient disagreement is imposed. Individual
components need not be bounded below in the class definition.

Let $\omega=(\pi_1,\ldots,\pi_K)$ consist of independent uniform
permutations of $\{1,\ldots,n\}$. If $t=en+j$ with $0\le e<K$ and
$0\le j<n$, write $i_t=\pi_{e+1}(j+1)$. We study the literal updates
\begin{equation}\label{eq:algorithm}
 x_{t+1}=x_t-\eta_t\nabla f_{i_t}(x_t),\qquad
 0<\eta_t\le\frac1L,
\end{equation}
and outputs
\begin{equation}\label{eq:output}
 \widehat x=\sum_{t=0}^T w_t x_t,
 \qquad w_t\ge0,\qquad \sum_{t=0}^T w_t=1.
\end{equation}
For the lower bound we allow $\eta_t$ and $w_t$ to be arbitrary
measurable functions of the instance, $\omega$, and an arbitrary private
seed independent of $\omega$. The constraints in
\eqref{eq:algorithm}--\eqref{eq:output} need hold only almost surely.
In particular, the selection may use the whole permutation bank before
the first update. This contains policies that see only the current
epoch's permutation and choose a step after observing the current
gradient. The choice of scalar step and convex output is unrestricted
within this model; the update direction and sampled order remain those
in \eqref{eq:algorithm}. Expectations of the nonnegative loss are
understood in the extended sense.

Define
\begin{align}
 A_0&=LD^2,& B_0&=\frac{LD^2}{T},\notag\\
 N_1&=\frac{\sigma D}{\sqrt T},&
 N_2&=\left(\frac{L\sigma^2D^4}{nK^2}\right)^{1/3},\label{eq:rates}\\
 N&=\min\{N_1,N_2\},&
 \Psi&=\min\{A_0,B_0+N\}.\notag
\end{align}

\begin{theorem}\label{thm:main}
For each $n\ge2$, $K\ge1$, $L>0$, and $D,\sigma\ge0$, there is one
instance in $\Class(n,L,D,\sigma)$, of dimension at most $n$, such that
every selection in \eqref{eq:algorithm}--\eqref{eq:output} satisfies
\begin{equation}\label{eq:lower}
 \E\bigl[f(\widehat x)-f(x_*)\bigr]\ge\frac{\Psi}{524288}.
\end{equation}
The instance is chosen before the seed space, seed law, and step/output
selection. Conversely, a parameter-dependent constant-step procedure,
using either the uniform pre-update average or the initial point,
satisfies the class-uniform upper bound $57\Psi$.
\end{theorem}

The dimension bound concerns the hard witness, not the definition of
$\Class$. For any fixed private-seed probability space, let $\Pol$
be the class of instance-dependent selections above, and let
$\Risk(I,P)$ denote their expected actual loss. The theorem gives
\begin{equation}\label{eq:minimax}
 \frac{\Psi}{524288}
 \le\sup_{I\in\Class}\inf_{P\in\Pol}\Risk(I,P)
 \le\inf_{P\in\Pol}\sup_{I\in\Class}\Risk(I,P)
 \le57\Psi.
\end{equation}
The lower is uniform even over the enlarged full-bank policy class.
The upper uses one procedure from its much smaller constant-step
subclass. Neither an upper for every policy nor equality of the two
minimax values is asserted.

\paragraph{Relation to previous results.}
Liu~\cite{liu2026} establishes the smooth-convex upper rate used here,
including a theorem for steps that vary by epoch. The powers of the
parameters in \eqref{eq:rates} are therefore inherited. Cha, Lee, and
Yun~\cite{cha2023} prove constant-step lower bounds for weighted
epoch-iterate outputs, as well as results for chosen permutations.
Rajput, Lee, and Papailiopoulos~\cite{rajput2022} give fixed-instance
constant-step obstructions and discuss variable-step lower bounds.
Thus fixed-instance constructions, weighted-output lower bounds, and
the cubic noise rate are prior ingredients. The cone witness below
itself has globally bounded gradient disagreement,
$\norm{\nabla f_i(x)-\nabla f(x)}=\sigma$; weakening that assumption
is not a contribution of the construction. The claim here is a
matching lower for adaptive scalar updates, all-state convex outputs,
and all $n,K$ and parameter regimes under a uniform smoothness bound
$L$. It addresses the optimized-rate part of the lower-bound question
in Liu's Section~5; it does not give the risk profile of every
prescribed schedule, match the separate nonuniform-smoothness
refinements, or settle every question posed there. It also does not
cover momentum, vector-valued
preconditioners, extra update directions, or outputs outside the convex
hull of the iterates.

\section{The inherited upper bound}

Theorem~1 of Liu~\cite{liu2026}, with all component smoothness bounds
set to $L$, implies that a constant $0<\eta\le1/(6L)$ and the
\emph{pre-update} average
\begin{equation}\label{eq:preaverage}
 \overline x=\frac1T\sum_{t=0}^{T-1}x_t
\end{equation}
satisfy
\begin{equation}\label{eq:liu}
 \E[f(\overline x)-f(x_*)]
 \le\frac{6D^2}{\eta T}
       +51\sigma^2\min\{\eta,nL\eta^2\}.
\end{equation}
This is precisely an admissible output in \eqref{eq:output}, with
$w_T=0$. In particular it is not an epoch-end average or a final-iterate
guarantee.

For $D,\sigma>0$, use either
\[
 \eta_1=\min\!\left\{\frac1{6L},\frac{D}{\sigma\sqrt T}\right\},
 \qquad
 \eta_2=\min\!\left\{\frac1{6L},
       \left(\frac{D^2}{nL\sigma^2T}\right)^{1/3}\right\}.
\]
Since the reciprocal of a minimum is a maximum, \eqref{eq:liu}
gives respectively $36B_0+57N_1$ and $36B_0+57N_2$. Choosing the
smaller noise branch gives at most $57(B_0+N)$. Smoothness of the mean
and $\nabla f(x_*)=0$ also give
$f(x_0)-f(x_*)\le A_0/2$. Choosing between this initial-point output
and \eqref{eq:preaverage}, using only the parameters, yields
\[
 \min\{A_0/2,57(B_0+N)\}\le57\Psi.
\]
When $\sigma=0<D$, use $\eta=1/(6L)$ and the same initial-point
option. When $D=0$, return $x_0$ while taking any legal positive steps.
These cases require no division by zero. This proves the upper part of
Theorem~\ref{thm:main} from the cited result. The remainder proves the
lower independently of \eqref{eq:liu}.

\section{A legal cone instance in every component count}

For now assume $D,\sigma>0$. Set $\theta=2^{-14}$ and choose
$0<\ell\le L$ with
\begin{equation}\label{eq:scale-condition}
 \sigma^2K\ge\ell^2D^2n.
\end{equation}
The final choice of $\ell$ will cover every parameter regime. Put
\begin{align}
 \mu&=\theta\min\!\left\{\ell,\frac{\sigma}{D\sqrt T},
       \left(\frac{\ell\sigma^2}{nK^2D^2}\right)^{1/3}\right\},
 &a&=\sqrt{\frac\mu\ell},\label{eq:mu}\\
 \beta&=\frac{\ell}{4(1+a^2)},&
 b&=\beta a,& H&=\mu+ba.\label{eq:coefficients}
\end{align}

\subsection{Fixed regular-simplex blocks}
For $n=2$ take a single block of size two. For $n\ge3$, write
$n=3q+j$, $j\in\{0,1,2\}$. If $j=0$, take $q$ blocks of size
three. If $j=1$ or $j=2$, take respectively one block of size four or
five and $q-1$ blocks of size three. This includes $n=3,4,5$ without
an exceptional construction. If $r$ is the number of blocks, then
\begin{equation}\label{eq:block-count}
 r\le n/3\quad(n\ge3),\qquad r\le n/2\quad(n\ge2).
\end{equation}
All blocks and label assignments are fixed before any permutations.

For a block of size $m$, use the zero-sum subspace of $\R^m$ and the
vectors
\begin{equation}\label{eq:simplex}
 \xi_i=\sqrt{\frac m{m-1}}
       \left(e_i-\frac1m\mathbf1\right),\qquad 1\le i\le m.
\end{equation}
They have norm one, sum zero, and satisfy
$\ip{\xi_i}{\xi_j}=-1/(m-1)$ for $i\ne j$. Different blocks occupy
orthogonal subspaces. Their direct sum $V$ has dimension $n-r$.
The same notation $\xi_i$ denotes the resulting fixed vectors indexed
by the original $n$ labels.

\subsection{Objective and actual gradients}
On the Euclidean product $\R\oplus V$, define the closed convex cone
and potential
\begin{align}
 \mathcal C_a&=\{(u,v):\norm v\le-au\},\notag\\
 \Phi(u,v)&=\frac\mu2u^2+\frac\ell8\norm v^2
                 +\frac\ell8\dist((u,v),\mathcal C_a)^2,\label{eq:potential}\\
 f_i(u,v)&=\Phi(u,v)+\sigma\ip{\xi_i}{v},
 &x_0&=(-D,0).\label{eq:components}
\end{align}
The mean is exactly $\Phi$, with minimizer $(0,0)$. We recall the
projection fact used to check global regularity. For a nonempty closed
convex set $C$, its projection $P_C$ satisfies
$\ip{x-P_Cx}{z-P_Cx}\le0$ for $z\in C$. Applying this twice shows
that $r(x)=x-P_Cx$ is $1$-Lipschitz. If
$g(x)=\norm{r(x)}^2/2$, comparison with the fixed competitors $P_Cx$
and $P_Cy$ gives
\[
 g(y)\le g(x)+\ip{r(x)}{y-x}+\tfrac12\norm{y-x}^2,\qquad
 g(y)\ge g(x)+\ip{r(y)}{y-x}-\tfrac12\norm{y-x}^2.
\]
It follows that $\nabla g=r$. Convexity follows by using a convex
combination of the two projected points as a feasible competitor.
Since $\mu\le\theta\ell<\ell/4$, the quadratic
part of \eqref{eq:potential} has gradient Lipschitz constant $\ell/4$.
The distance term contributes another $\ell/4$. Thus all components
are convex and globally $\ell/2$-smooth, hence $L$-smooth. Their
gradients at zero are $(0,\sigma\xi_i)$, so the variance in
\eqref{eq:parameters} is exactly $\sigma^2$. The initial distance is
$D$, and the dimension is $n-r+1\le n$. In fact these components are
bounded below, although the public class does not require it.

Write a reachable point as $(-q,v)$, with $q>0$ and $R=\norm v$.
If $R\le aq$, it is in the cone. Otherwise its projection is
\begin{equation}\label{eq:projection}
 \left(-\widehat q,\frac{a\widehat q}{R}v\right),
 \qquad \widehat q=\frac{q+aR}{1+a^2}.
\end{equation}
For completeness, at a fixed candidate slow coordinate $-s$, $s\ge0$,
the nearest fast point is the projection of $v$ onto the ball of radius
$as$. Minimizing $(s-q)^2+(R-as)_+^2$ gives
\eqref{eq:projection} in the outside case; here
$0<\widehat q<R/a$. This also verifies the formula at the boundary
by continuity.

Consequently the \emph{actual} update for current offset $\xi$ is
\begin{align}
 q'&=(1-\mu\eta)q+b\eta(R-aq)_+,
 &v'&=(1-d\eta)v-\sigma\eta\xi,\label{eq:dynamics}\\
 d&=\frac\ell4+\beta\frac{(R-aq)_+}{R}.\notag
\end{align}
The fraction is defined to be zero at $R=0$. In particular,
\begin{equation}\label{eq:retention-basic}
 \frac\ell4\le d\le\frac\ell2,\qquad
 \rho:=1-d\eta\ge\frac12,\qquad
 q'\ge(1-\mu\eta)q>0.
\end{equation}
These inequalities use the original cap $\eta\le1/L\le1/\ell$.
Starting from $q_0=D$, they justify the reachable-state formulas at
every update, for every legal selection.

\begin{lemma}[Radius and energy]\label{lem:energy}
Every actual path satisfies $R_t\le4\sigma/\ell$. On every interval
$J=[s,t)\subseteq[0,T)$, with
$S_J=\sum_{k\in J}\eta_k$ and $W_J=\sum_{k\in J}\eta_kR_k$,
\begin{equation}\label{eq:energy}
 \sum_{k\in J}\eta_k^2
 \le\frac{R_t^2-R_s^2}{\sigma^2}+\frac{6W_J}{\sigma}.
\end{equation}
\end{lemma}
\begin{proof}
The norm update gives
$R'\le(1-\ell\eta/4)R+\sigma\eta$. The asserted radius follows
by induction from $R_0=0$. Moreover,
\begin{align*}
 R'^2-R^2
 &=-(2d\eta-d^2\eta^2)R^2
       -2\rho\sigma\eta\ip v\xi+\sigma^2\eta^2\\
 &\ge\sigma^2\eta^2-\ell\eta R^2-2\sigma\eta R
 \ge\sigma^2\eta^2-6\sigma\eta R.
\end{align*}
Sum this inequality over the actual interval. No conditional
expectation or independence of the step and label is used.
\end{proof}

\section{Dual tables from the sampled permutations}

Fix the entire permutation bank. In one epoch, a block of size $m$
occupies positions $p_1<\cdots<p_m$ in $\{0,\ldots,n-1\}$. Choose
a largest consecutive gap $g=p_{k+1}-p_k$, with deterministic
tie-breaking, and let $E$ be its $k$ earlier labels. Set
$z=\sum_{i\in E}\xi_i$. The fixed simplex identities give
\[
 \ip z{\xi_i}=\begin{cases}(m-k)/(m-1),&i\in E,\\
                         -k/(m-1),&i\notin E,
             \end{cases}\qquad
 \norm z^2=\frac{k(m-k)}{m-1}.
\]
Use dual vectors
\begin{equation}\label{eq:dual-endpoints}
 \phi^-=\frac{m-1}{m-k}z,\qquad
 \phi^+=-\frac{m-1}{k}z.
\end{equation}
Keep $\phi^-$ up to $p_k$, interpolate linearly to $\phi^+$ across
the gap, and keep $\phi^+$ thereafter.
There is no hit from this block inside the chosen gap. Hence the dual
has inner product one with each actual offset at its hit position.
Also
\[
 \norm{\phi^-}^2=\frac{(m-1)k}{m-k},\quad
 \norm{\phi^+}^2=\frac{(m-1)(m-k)}k,\quad
 \norm{\phi^+-\phi^-}
   =m\sqrt{\frac{m-1}{k(m-k)}}\le m.
\]
For $2\le m\le5$, the endpoint norms are at most $4$ and the
change is at most $5$.

Taking the orthogonal direct sum over all blocks gives a dual
$\phi_t$ at every update position such that
\begin{equation}\label{eq:dual-bounds}
 \ip{\phi_t}{\xi_{i_t}}=1,\qquad
 \norm{\phi_t}\le4\sqrt r,\qquad
 |\ip{\phi_t}{\xi_i}|\le4\quad\hbox{for every original label }i.
\end{equation}
The last bound holds even when the label and the dual time differ.
This fact is needed when paying for resets.

Write $g_{e,j}$ for the selected gap of block $j$ in epoch $e$, and
define
\begin{equation}\label{eq:geometry}
 B_e=\sum_{j=1}^r\frac1{g_{e,j}},\qquad
 G_A(\omega)=A\left[8K\sqrt r+
             5\sqrt{n-1}\sum_{e=1}^K\sqrt{B_e}\right].
\end{equation}
Within an epoch, the squared variation of block $j$ is at most
$25/g_{e,j}$. Orthogonality and Cauchy--Schwarz therefore give
\begin{equation}\label{eq:variation}
 \sum_{\text{internal }t}\norm{\phi_{t+1}-\phi_t}^2\le25B_e,\qquad
 \sum_{\text{internal }t}\norm{\phi_{t+1}-\phi_t}
       \le5\sqrt{(n-1)B_e}.
\end{equation}
An epoch-boundary jump has norm at most $8\sqrt r$. Thus for any
interval $[s,t)$ and any vectors $z_s,\ldots,z_t$ bounded in norm by
$A$, summation by parts gives
\begin{equation}\label{eq:sbp}
 \left|\sum_{k=s}^{t-1}\ip{\phi_k}{z_k-z_{k+1}}\right|
 \le G_A(\omega).
\end{equation}
Indeed the two endpoints cost at most $8A\sqrt r$, at most $K-1$
boundary jumps cost $8A(K-1)\sqrt r$, and the internal variations
are bounded by \eqref{eq:variation}. This is simultaneous over all
intervals and all such vector sequences.

\begin{lemma}[Expected geometry]\label{lem:geometry}
For every $n\ge2$, $\E B_e\le4$, and for deterministic $A\ge0$,
\begin{equation}\label{eq:meanG}
 \E G_A\le16AK\sqrt n.
\end{equation}
\end{lemma}
\begin{proof}
For $n\ge3$, fix any three labels of a block. Their positions form
a uniform three-element subset: each position subset is realized by
$3!(n-3)!$ of the $n!$ permutations. If its span is $s$, there are
$(n-s)(s-1)$ such subsets, so
\[
 \E\frac1s
 =\binom{n}{3}^{-1}\sum_{s=2}^{n-1}\frac{(n-s)(s-1)}s
 \le\binom{n}{3}^{-1}\frac{(n-1)(n-2)}2=\frac3n.
\]
The outer span of the whole block is at least $s$, and its largest
gap is at least its outer span divided by $m-1\le4$. Thus
$\E(1/g_{e,j})\le12/n$. By \eqref{eq:block-count},
$\E B_e\le4$. For $n=2$ there is one gap, equal to one.
Finally Jensen gives $\E\sqrt{B_e}\le2$, and
\[
 \E G_A\le AK\bigl(8\sqrt{n/2}+10\sqrt{n-1}\bigr)
 <16AK\sqrt n
\]
when $A>0$. The non-strict version also covers $A=0$.
Only the marginal positions of fixed labels were used. The duals may
be correlated with every action of the policy.
\end{proof}

\section{Paying for excursions of the actual trajectory}

The next estimate is pathwise. Fix a bank, a legal realized trajectory,
an interval $J=[s,t)$, and deterministic $h,A>0$ with $A>\sigma h$.
Keep update $k$ if
\begin{equation}\label{eq:keep}
 \eta_k<h,\qquad R_k\le A,\qquad R_{k+1}\le A,
\end{equation}
and remove the other updates. Subscripts $\mathrm{keep}$ and
$\mathrm{rem}$ denote the corresponding partial sums.

\begin{lemma}[Paid auxiliary resets]\label{lem:paid}
For every such interval,
\begin{equation}\label{eq:paid}
 S_{\mathrm{keep}}\le\frac{G_A}{\sigma}+4S_{\mathrm{rem}}
       +\frac{2\ell\sqrt r}{\sigma}W_J
       +\frac{4\sqrt r}{\sigma}R_s.
\end{equation}
\end{lemma}
\begin{proof}
Construct an auxiliary sequence solely for accounting. Set $z_s=v_s$
if $R_s\le A$, and $z_s=0$ otherwise. Copy the actual path on kept
updates. During each maximal removed run freeze the auxiliary sequence
until the run's final update. If the actual state at the run's endpoint
has radius at most $A$, reset there to that state. Otherwise leave it
frozen. Such an unsafe endpoint can occur only on a terminal removed
run: a following kept update would require that same state to be safe.
Thus the construction agrees with $v$ at both endpoints of every kept
update and remains within radius $A$ at every time.

For a removed run ending safely, the reset displacement is the sum of
the actual fast increments along the run. The sole possible additional
term is $v_s$, when the first run starts at an unsafe initial state.
Test the reset against the dual at its final update. For every update
inside that run, \eqref{eq:dynamics} and the all-label bound in
\eqref{eq:dual-bounds} give
\[
 |\ip{\phi_j}{v_{k+1}-v_k}|
 \le4\sigma\eta_k+2\ell\sqrt r\,\eta_kR_k
 \quad\hbox{for every test time }j.
\]
Consequently all reset charges together are at most
\begin{equation}\label{eq:reset-payment}
 4\sigma S_{\mathrm{rem}}+2\ell\sqrt r\,W_{\mathrm{rem}}
       +4\sqrt r\,R_s.
\end{equation}
Every removed actual increment is charged at most once. On a kept
update, \eqref{eq:dual-bounds} yields the exact identity
\[
 \sigma\eta_k=\ip{\phi_k}{z_k-z_{k+1}}
                  -d_k\eta_k\ip{\phi_k}{v_k}.
\]
Sum over kept updates, insert the fully paid resets, and use
\eqref{eq:sbp}. The kept drift terms cost at most
$2\ell\sqrt r\,W_{\mathrm{keep}}$. Combining with
\eqref{eq:reset-payment} and dividing by $\sigma$ proves
\eqref{eq:paid}.

This auxiliary sequence is not a modified algorithm. In particular no
actual state is reset or projected, no excursion is discarded without
payment, and a removed run may cross epoch boundaries. The argument
works for intervals and keep/remove choices selected after the full
trajectory has been seen.
\end{proof}

\subsection{An unrestricted interval estimate}
Set $C=HD$ and choose
\begin{equation}\label{eq:thresholds}
 b_0=b-\frac{2C\ell\sqrt r}{\sigma},\qquad
 h=\frac{120C}{b_0\sigma},\qquad A=\frac{130C}{b_0}.
\end{equation}
We first check that these expressions are positive. From
\eqref{eq:scale-condition} and the linear cap in \eqref{eq:mu},
\[
 a^2\le\theta,\qquad
 a^2\le\frac{\theta\sigma^2}{\ell^2D^2n},\qquad
 \frac Hb=a(5+4a^2).
\]
Together with $r\le n/2$, these give
\[
 \frac{2C\ell\sqrt r}{\sigma b}
 \le2(5+4\theta)\sqrt{\theta/2}<\frac1{16}.
\]
For the last inequality, use $5+4\theta<21/4$ and
$\sqrt\theta=1/128$, obtaining $21/(256\sqrt2)<1/16$.
In particular,
\begin{equation}\label{eq:b0}
 \frac{15b}{16}\le b_0\le b,\qquad
 A-\sigma h=\frac{10C}{b_0}>0.
\end{equation}

\begin{lemma}[All intervals]\label{lem:interval}
With \eqref{eq:thresholds}, every actual interval $J=[s,t)$ satisfies
\begin{equation}\label{eq:interval}
 C S_J-bW_J
 \le\frac{C G_A}{\sigma}
       +\frac{4C\sqrt r}{\sigma}R_s
       +\frac{b_0}{24\sigma}R_t^2.
\end{equation}
\end{lemma}
\begin{proof}
Add $S_{\mathrm{rem}}$ to \eqref{eq:paid}, multiply by $C$, and
subtract $bW_J$. The result is
\begin{equation}\label{eq:preinterval}
 CS_J-bW_J\le\frac{CG_A}{\sigma}
       +\frac{4C\sqrt r}{\sigma}R_s
       +5CS_{\mathrm{rem}}-b_0W_J.
\end{equation}
Every large update, meaning $\eta_k\ge h$, is removed. By
Lemma~\ref{lem:energy},
\begin{align}
 5CS_{\mathrm{large}}
 &\le\frac{5C}{h}\sum_{k\in J}\eta_k^2\notag\\
 &\le\frac{b_0}{24\sigma}(R_t^2-R_s^2)
                  +\frac{b_0}4 W_J.\label{eq:large}
\end{align}
This is valid even if $h>1/L$ and there are no large updates: the
right-hand bracket before multiplying by $5C/h$ is nonnegative by
\eqref{eq:energy}.

A small removed update has at least one unsafe endpoint. Since
$R_{k+1}\le R_k+\sigma\eta_k$, its pre-update radius satisfies
$R_k>A-\sigma h=10C/b_0$. Therefore
\begin{equation}\label{eq:small}
 5CS_{\mathrm{small,rem}}\le\frac{b_0}2W_{\mathrm{small,rem}}.
\end{equation}
Insert \eqref{eq:large}--\eqref{eq:small} into
\eqref{eq:preinterval}. The remaining weighted-radius terms are
nonpositive, as is the starting-energy term. Dropping them proves
\eqref{eq:interval}.
\end{proof}

\section{Endpoint trimming and a bank-only certificate}

For an actual path define its maximal interval deficit
\begin{equation}\label{eq:path-deficit}
 Z_{\mathrm{path}}=\max\!\left\{0,
       \max_{0\le s<t\le T}\sum_{k=s}^{t-1}\eta_k(C-bR_k)\right\}.
\end{equation}
The global radius bound alone is too costly in the endpoint terms of
\eqref{eq:interval}. The following finite selection argument retains
small endpoints instead.

\begin{lemma}[Endpoint trim]\label{lem:trim}
Every legal path for a fixed bank satisfies
\begin{equation}\label{eq:certificate}
 Z_{\mathrm{path}}\le\mathcal Z(\omega):=
 \frac CL+\frac{C G_A(\omega)}{\sigma}
       +\frac{4C^2\sqrt r}{\sigma b}+\frac{C^2}{24\sigma b}.
\end{equation}
The certificate $\mathcal Z$ is nonnegative and depends only on the
parameters and permutation bank, not on the steps or output.
\end{lemma}
\begin{proof}
If $Z_{\mathrm{path}}>0$, choose a shortest interval attaining it.
Its first and last increments are strictly positive: a negative
endpoint could be deleted to increase the maximum, while a zero
endpoint could be deleted without changing it. As all steps are
positive, $R_s<C/b$ and $R_{t-1}<C/b$.

Delete the final update, whose contribution is at most $C/L$. If the
remainder is empty, this already bounds the maximum. Otherwise apply
\eqref{eq:interval} to $[s,t-1)$. Its initial and final state radii
are both less than $C/b$. Using $b_0\le b$ gives the three remaining
terms of \eqref{eq:certificate}. The case $Z_{\mathrm{path}}=0$ is
immediate. This argument conditions on no selected endpoint event.
\end{proof}

\begin{lemma}[Mean certificate]\label{lem:mean-certificate}
The construction satisfies
\begin{equation}\label{eq:small-certificate}
 \E\mathcal Z<\frac D{128}.
\end{equation}
\end{lemma}
\begin{proof}
By \eqref{eq:b0}, $A\le2080C/(15b)$. Use
Lemma~\ref{lem:geometry}, $\sqrt r\le K\sqrt n$, and
$1\le K\sqrt n$ in \eqref{eq:certificate} to obtain
\begin{align}
 \frac{\E\mathcal Z}{D}
 &\le\frac HL+
   \left(\frac{33280}{15}+4+\frac1{24}\right)
             \frac{H^2DK\sqrt n}{\sigma b}\notag\\
 &=\frac HL+\frac{53345}{24}
             \frac{H^2DK\sqrt n}{\sigma b}.\label{eq:certificate-mean}
\end{align}
The coefficients in \eqref{eq:coefficients} obey
\begin{equation}\label{eq:H-b}
 H\le\frac54\mu,\qquad
 \frac{H^2}{b}\le\frac{25}{4}(1+\theta)
                     \frac{\mu^{3/2}}{\sqrt\ell}
                  <7\frac{\mu^{3/2}}{\sqrt\ell}.
\end{equation}
The linear and cubic caps defining $\mu$ therefore imply
\[
 \frac{\E\mathcal Z}{D}
 <\frac54\theta+7\theta^{3/2}\frac{53345}{24}<\frac1{128}.
\]
For a finite check, $53345/24<2240$, so the middle expression is
less than $(160+15680)/2^{21}<1/128$.
\end{proof}

\section{Retention of every state and adaptive output}

For a pre-state with $q_k\le D$, \eqref{eq:dynamics} gives
\begin{equation}\label{eq:slow-drop}
 q_{k+1}\ge q_k-C\eta_k+b\eta_kR_k.
\end{equation}
Indeed $(R-aq)_+\ge R-aq$ and $C=(\mu+ba)D$.

\begin{proposition}\label{prop:cone-risk}
For the one cone instance above, every allowed selection satisfies
\begin{equation}\label{eq:cone-risk}
 \E\Phi(\widehat x)\ge\frac{\mu D^2}{16}.
\end{equation}
\end{proposition}
\begin{proof}
Suppose an actual state first falls below $D/2$ at time $t$. If there
is no earlier state with $q_j>D$, put $s=0$. Otherwise let $j<t$ be
the last such time and put $s=j+1$. In this second case,
\eqref{eq:retention-basic} implies
$q_s>(1-\mu/L)D$; in the first case $q_s=D$. Since
$\mu/L\le\theta<1/4$, we have $s<t$. All pre-states in $[s,t)$
satisfy $q_k\le D$. Summing \eqref{eq:slow-drop} shows
\[
 Z_{\mathrm{path}}\ge\sum_{k=s}^{t-1}\eta_k(C-bR_k)
          \ge q_s-q_t>(1/2-\mu/L)D>D/4.
\]
By Lemma~\ref{lem:trim}, on every bank with
$\mathcal Z(\omega)\le D/4$ \emph{all} legal selections have
$q_t\ge D/2$ for every $0\le t\le T$. This event is independent
of the choice of policy in the sense that its definition involves
only the bank; no statistical independence of actions is needed.
Markov's inequality and \eqref{eq:small-certificate} give
\[
 \Prob\{\mathcal Z>D/4\}<1/32.
\]
Every scalar convex output on the complementary event has slow
coordinate $-\widehat q$ with $\widehat q\ge D/2$, even when the
weights use the entire trajectory and seed. Since
$\Phi(u,v)\ge\mu u^2/2$ and $\Phi\ge0$ everywhere,
\[
 \E\Phi(\widehat x)\ge\frac{31}{32}\frac{\mu D^2}{8}
                         \ge\frac{\mu D^2}{16}.
\]
For product-almost-everywhere legal selectors, the same argument
applies outside their null set. Continuity of the actual gradients
and measurable selectors makes the finitely many states and the loss
measurable. Nonnegative integration requires no additional
integrability assumption. The instance and its good-bank event are
unchanged for every independent private seed law.
\end{proof}

\section{All regimes and the single-instance choice}

For arbitrary $D,\sigma>0$, choose
\begin{equation}\label{eq:ell}
 \ell=\min\!\left\{L,\frac{\sigma\sqrt K}{D\sqrt n}\right\}.
\end{equation}
This is a smoothness scale of the constructed instance; the algorithm
still obeys the original cap $1/L$. It ensures
\eqref{eq:scale-condition}. Let $M=\min\{A_0,N_1,N_2\}$.
If $\ell=L$, \eqref{eq:mu} immediately gives $\mu D^2=\theta M$.
If $\ell<L$, put $s_0=\sigma/(D\sqrt T)$. Then $\ell=Ks_0$,
and the cubic term in \eqref{eq:mu} is exactly $s_0$. Thus
$\mu=\theta s_0$. Moreover
$\sigma^2K<L^2D^2n$ implies $N_1<A_0$ and $N_1<N_2$.
We conclude in both cases that
\begin{equation}\label{eq:mu-rate}
 \mu D^2=\theta M,\qquad
 \E\Phi(\widehat x)\ge\frac{M}{262144}.
\end{equation}

The other candidate is the identical scalar quadratic
\begin{equation}\label{eq:quadratic}
 f_i(x)=\frac\lambda2x^2,\qquad
 \lambda=\frac{L}{2T},\qquad x_0=D,\qquad x_*=0.
\end{equation}
It is legal for every $\sigma\ge0$, since its optimum variance is
zero. Along every allowed realized schedule,
\[
 x_t=D\prod_{j=0}^{t-1}(1-\lambda\eta_j)
 \ge D\left(1-\sum_{j=0}^{t-1}\lambda\eta_j\right)\ge D/2.
\]
Here every factor lies in $[0,1]$, and the product bound follows
inductively. Any scalar convex output therefore satisfies
\begin{equation}\label{eq:quadratic-risk}
 f(\widehat x)-f(0)\ge\frac{\lambda D^2}{8}=\frac{B_0}{16}.
\end{equation}

Choose between \eqref{eq:components} and \eqref{eq:quadratic} by
comparing the two proved lower bounds, before any policy or randomness
is chosen. Since $B_0+M\ge\Psi$,
\[
 \max\!\left\{\frac{M}{262144},\frac{B_0}{16}\right\}
 \ge\frac{B_0+M}{524288}\ge\frac\Psi{524288}.
\]
This is one parameter-only witness, not an instance selected after
observing the strategy. For $\sigma=0<D$, use only the quadratic;
then $\Psi=B_0$. For $D=0$, $\Psi=0$, the quadratic is again legal,
and returning $x_0$ achieves zero. These branches avoid all undefined
cone parameters. The proof of Theorem~\ref{thm:main} is complete.

\section{Verification and scope}

The lower theorem has also been expressed in Lean~4 using Mathlib.
Its final statement is
\texttt{Gate149.FinalTheorem.}\allowbreak
\texttt{one\_fixed\_instance\_fresh\_epoch\_lower}.
It constructs an actual standard-Euclidean instance of dimension at
most $n$ before quantifying over the seed space, probability law, and
measurable full-bank selectors. The formal development includes the
globally convex smooth objective, its actual gradient recursion,
fixed-label simplex geometry, paid resets, interval trimming, parameter
tuning, and the joint probability law. In particular, the finite uniform
permutation-bank law is proved equal to the product of fresh uniform
epoch laws. The final lower statement has been accepted by the Lean
kernel, with its axiom report containing only \texttt{propext},
\texttt{Classical.choice}, and \texttt{Quot.sound}.

The formal lower and the inherited upper have different validation
scopes. The proof of Liu's published upper theorem is not included in
that development. Its minimax sandwich is an explicitly conditional
formal corollary given the upper strategy, while Section~2 supplies
the mathematical citation and tuning. No claim about priority or
external assessment follows from kernel acceptance.

The substantive lower-bound mechanism is the simultaneous pathwise
charge \eqref{eq:paid}--\eqref{eq:interval}, followed by endpoint
trimming. It prevents full-order information, adaptive scalar steps,
and trajectory-selected convex averaging from improving the optimized
worst-case rate on this class. This is a statement about the specified
scalar-gradient model, not a limitation on all finite-sum optimization
methods or a prediction of practical training performance.

\begin{thebibliography}{9}
\small
\setlength{\itemsep}{4pt}

\bibitem{liu2026}
Zijian Liu.
\newblock Random reshuffling dominates stochastic gradient descent.
\newblock In \emph{Proceedings of the 39th Conference on Learning Theory},
volume 336 of \emph{Proceedings of Machine Learning Research},
pages 4859--4882, 2026.
\newblock \url{https://proceedings.mlr.press/v336/liu26d.html}.

\bibitem{cha2023}
Jaeyoung Cha, Jaewook Lee, and Chulhee Yun.
\newblock Tighter lower bounds for shuffling SGD: Random permutations
and beyond.
\newblock In \emph{Proceedings of the 40th International Conference on
Machine Learning}, volume 202 of \emph{Proceedings of Machine Learning
Research}, pages 3855--3912, 2023.
\newblock \url{https://proceedings.mlr.press/v202/cha23a.html}.

\bibitem{rajput2022}
Shashank Rajput, Kangwook Lee, and Dimitris Papailiopoulos.
\newblock Permutation-based SGD: Is random optimal?
\newblock In \emph{International Conference on Learning Representations},
2022.
\newblock \url{https://openreview.net/forum?id=YiBa9HKTyXE}.
\newblock Full version: \url{https://arxiv.org/abs/2102.09718v2}.

\end{thebibliography}
\end{document}
